
\documentclass[10pt,a4paper]{article}

% ---------- Page setup ----------
\usepackage[
  top=17mm,
  bottom=17mm,
  left=19mm,
  right=19mm
]{geometry}

\usepackage[T1]{fontenc}
\usepackage[utf8]{inputenc}
\usepackage{lmodern}
\usepackage{microtype}
\usepackage{enumitem}
\usepackage{titlesec}
\usepackage{xcolor}

% ---------- Color palette ----------
\definecolor{ink}{HTML}{193A4D}
\definecolor{muted}{HTML}{52616B}
\definecolor{rulegray}{HTML}{D4DEE3}

% ---------- Typography ----------
\setlength{\parindent}{0pt}
\setlength{\parskip}{3pt}
\setlength{\emergencystretch}{1.5em}
\linespread{1.055}

\raggedbottom
\pagestyle{empty}

% ---------- Section formatting ----------
\titleformat{\section}
  {\color{ink}\bfseries\normalsize}
  {}{0pt}{}
  [\vspace{2pt}\color{rulegray}\titlerule]

\titlespacing*{\section}
  {0pt}{9pt}{4pt}

% ---------- List formatting ----------
\setlist[itemize]{
  leftmargin=1.25em,
  itemsep=4pt,
  topsep=3pt,
  parsep=0pt,
  partopsep=0pt
}

\begin{document}

% ==================================================
% Title
% ==================================================

\begin{center}
  {\fontsize{15}{18}\selectfont
   \bfseries\color{ink}
   DecisionRL: Enhancing Decision Calibration\\
   through Reinforcement Learning\par}

  \vspace{4pt}

  {\small\color{muted}
   Reinforcement Learning Course Project Proposal}

  \vspace{5pt}

  {\small
   \textbf{Team Members:}
    \quad\textbullet\quad
   Taihua Chen (260010020)
   \quad\textbullet\quad
   Yanyun LIU (260010215)}
\end{center}

\vspace{1pt}

% ==================================================
% Research Question
% ==================================================

\section*{Research Question}

As large language models evolve from
question answerers to decision makers,
efficient and reliable decision-making becomes
increasingly important.
While System 2-style models emphasize deliberate,
multi-step reasoning, System 1-style models aim to
provide fast, intuitive judgments.
Jev introduces a complementary paradigm of
structured probabilistic decisions, producing
explicit probability distributions over candidate
actions instead of lengthy reasoning trajectories.
Inspired by this paradigm, our project explores how
\textbf{Reinforcement Learning for Calibrated
Decisions (RLCD)} can enhance decision accuracy,
probability calibration, and reliability through
a reproducible training framework.

% ==================================================
% Related Work
% ==================================================

\section*{Related Work}

% \textbf{Jev} introduces System 1-style probabilistic
% decision-making and RLCD as an alternative to
% autoregressive reasoning.
% However, its model weights, training code, and full
% training recipe remain proprietary, limiting
% independent reproduction.
% \textbf{Laya} provides an open-source implementation
% of Jev-compatible typed decisions and RLCD,
% offering an accessible foundation for further research.
% \textbf{Sev} investigates RLCD optimization through
% independent reproduction and analysis.
% \textbf{Kev} and \textbf{Strands Decider} extend the
% open-source decision-model ecosystem through
% alternative architectures and training pipelines. Despite these advances, RLCD optimization remains
% relatively underexplored, particularly in reward
% design and policy optimization.
% Our project aims to develop a reproducible RLCD
% training framework and investigate improvements
% to its optimization.

\textbf{Jev} introduces System 1-style probabilistic
decision-making and RLCD, while
\textbf{Microsoft-Decision-1} demonstrates the potential
of efficient decision models.
However, their proprietary implementations and
limited disclosure of training methodologies
restrict reproducibility.
\textbf{Clef} provides open model weights but
does not offer a fully reproducible training pipeline. Meanwhile, \textbf{Laya} and \textbf{Kev} provide
open-source implementations of typed-decision models.
Laya introduces a practical RLCD training approach,
whereas Kev primarily adopts supervised learning
and post-hoc calibration.
However, RLCD optimization remains relatively
underexplored, particularly in reward design
and policy optimization.
Our project aims to develop a reproducible RLCD
training framework and explore improvements
in reinforcement learning strategies.


% ==================================================
% Planned Approach
% ==================================================


\section*{Planned Approach}

We plan to develop \textbf{DecisionRL} by exploring
decision model design and reinforcement learning
strategies for reliable and calibrated decision-making.
This development will involve two main components:
data construction and curation, and model training
with RLCD optimization.

\begin{itemize}

  \item \textbf{Data Construction and Curation.}
  We will use all official \texttt{typed-decisions}
  training examples and collect additional
  open-source decision datasets across multiple
  domains.
  Our data preparation will cover two types of
  supervision: hard labels for SFT and probabilistic
  calibration labels for RLCD.
  Since reliable probabilistic labels are scarce,
  we will synthesize diverse decision scenarios
  and construct calibration labels from
  teacher-model logit distributions or aggregated
  predictions through repeated sampling and voting.

  \item \textbf{Model Training and RLCD Optimization.}
  Using \textbf{EuroBERT} as our base checkpoint,
  we will adopt a two-stage training framework.
  First, we will perform supervised fine-tuning
  (SFT) on high-quality decision datasets to establish
  general decision-making capabilities.
  We will then apply RLCD to optimize decision
  probabilities using verified probabilistic
  supervision.
  We also aim to improve Laya's original RLCD
  algorithm by exploring enhanced reinforcement
  learning optimization strategies for more
  reliable and calibrated decisions.

\end{itemize}


% ==================================================
% Evaluation Plan
% ==================================================

\section*{Evaluation Plan}

We will evaluate DecisionRL on hard-label benchmarks
(e.g. Banking77) and soft-label benchmarks
(e.g. Typed Decisions), comparing it with
representative decision models (e.g. Jev).
The evaluation will use \textbf{Accuracy},
\textbf{Macro F1 Score},
\textbf{Expected Calibration Error},
and \textbf{Brier Score}
to assess the foundational decision-making
capabilities established through SFT, with particular
emphasis on analyzing how RLCD and our proposed
improvements affect probability calibration
and decision-making behavior.

% ==================================================
% Resources and Responsibilities
% ==================================================

\section*{Required Resources and Division of Responsibilities}

\textbf{Resources:}
No additional equipment or specialized resources
are required.

\textbf{Preliminary Division of Responsibilities:}
Both team members will jointly participate in project
discussions, research planning, and RLCD methodological
design.
Yanyun LIU will primarily contribute to data
collection, synthesis, and quality control,
while Taihua Chen will focus more on model and algorithm optimization.

\end{document}
